%% kapa: exam=2 round=26 year=2015 period=3 subject=compensation_law subject_ko=감정평가 및 보상법규 seq=2 number=2 points=30
%% source: Q-Net 감정평가사 기출문제 (cst00309) artlSeq=5149601 file=extracted/1008/5149601/2015년 제26회 감정평가사 2차시험 기출문제/3교시(감정평가 및 보상법규).pdf
%% conversion: kordoc 4.24.51 -> tools/gen2.py (auto); review=proofread-gpt-sol; flags=none
\begin{longquestion}{【문제 2】}
B시에 거주하는 甲은 2005년 5월 자신의 토지 위에 주거용 건축물을 신축하였다. 그런데 甲은 건축허가 요건을 충족하지 못하여 행정기관의 허가 없이 건축하였다. 甲은 위 건축물에 입주하지 않았으나, 친척인 乙이 자신에게 임대해 달라고 요청하여 이를 허락하였다. 乙은 필요시 언제든 건물을 비워주겠으며, 공익사업시행으로 보상의 문제가 발생할 때에는 어떠한 보상도 받지 않겠다는 내용의 각서를 작성하여 임대차계약서에 첨부하였다. 乙은 2006년 2월 위 건축물에 입주하였는데, 당시부터 건축물의 일부를 임의로 용도변경하여 일반음식점으로 사용하여 왔다. 甲의 위 토지와 건축물은 2015년 5월 14일 국토교통부장관이 한 사업인정 고시에 따라서 공익사업시행지구에 편입되었다. 甲은 이 사실을 알고 동년 6월에 위 건축물을 증축하여 방의 개수를 2개 더 늘려 자신의 가족과 함께 입주하였다. 다음 물음에 답하시오. (30점) \par
1) 위 甲의 건축물은 「공익사업을 위한 토지 등의 취득 및 보상에 관한 법률」에 따른 손실보상의 대상이 되는지, 만일 된다면 어느 범위에서 보상이 이루어져야 하는지 설명하시오. (10점) \par
2) 甲과 乙은 주거이전비 지급 대상자에 포함되는지 여부를 그 지급 요건에 따라서 각각 설명하시오. (20점) \par
\end{longquestion}
